\question[5]
\begin{enumerate}[label=\alph*)]
\item Nenne die Komponenten eines Von-Neumann-Rechners und beschreibe jeweils kurz ihre Aufgabe.
\item Erkläre den Unterschied zwischen dem Adressbus und dem Datenbus.
\end{enumerate}

\begin{solutionbox}{8cm}
\begin{enumerate}[label=\alph*)]
\item
\begin{itemize}
\item \textbf{Rechenwerk} (ALU): führt Rechenoperationen und logische Verknüpfungen aus.
\item \textbf{Steuerwerk}: holt die Befehle nacheinander aus dem Speicher, entschlüsselt sie
und steuert die anderen Komponenten so, dass der Befehl ausgeführt wird.
\item \textbf{Speicherwerk} (Arbeitsspeicher): speichert Programme und Daten.
\item \textbf{Ein-/Ausgabewerk}: verbindet den Rechner mit Ein- und Ausgabegeräten.
\item \textbf{Bussystem}: verbindet die Komponenten (Adressbus, Datenbus, Steuerbus).
\item \textbf{Register}: schnelle Speicher im Prozessor, z.\,B. Akkumulator, Befehlsregister, Befehlszähler.
\end{itemize}
Rechenwerk und Steuerwerk bilden zusammen den Prozessor (CPU).
\item Über den \textbf{Adressbus} wird die Adresse übertragen, also welche Speicherzelle gelesen
oder beschrieben wird. Über den \textbf{Datenbus} wird der Inhalt übertragen, also die Daten
oder Befehle selbst.
\end{enumerate}
\end{solutionbox}
